\clearpage\thispagestyle{empty}
\pdfbookmark[0]{Sources for the review of Rossel’s second volume}{rossel1809-sources}
\begin{center}\textbf{Sources for the review of Rossel’s second volume}\end{center}
\noindent This AI-translated review follows Carl Friedrich Gauss, \href{https://gdz.sub.uni-goettingen.de/id/PPN236007467}{\emph{Werke}, vol.~VI} (Göttingen, 1874), pp.~520–522. The review first appeared unsigned under Paris in the \emph{Göttingische gelehrte Anzeigen}, no.~167, 21 October 1809, \href{https://images.p.sub.uni-goettingen.de/iiif/image/gdz:PPN319721507_1809:00001661/full/1650,/0/default.jpg}{pp.~1657}–\href{https://images.p.sub.uni-goettingen.de/iiif/image/gdz:PPN319721507_1809:00001667/full/1650,/0/default.jpg}{1663}. It discusses the second volume of the voyage of d’Entrecasteaux, sent in search of La Pérouse. The French title says it was published by order of the Emperor and King and prepared by M. de Rossel, a former naval captain. The \href{https://catalogue.bnf.fr/ark:/12148/cb43383186m}{Bibliothèque nationale de France catalogue} records Rossel as the scientific editor.
\par\smallskip\noindent The collected edition places a period after \emph{ausmacht} and has an apostrophe-like mark before \emph{den}. The first printing uses a comma there and has no such mark. It uses double-prime marks in the oscillation timings where the collected edition uses a raised \emph{s}. The translations follow the collected readings.
\par\smallskip\noindent In section V Gauss quotes Rossel in French.
\begin{quote}
Tout ce qui peut faciliter une observation aussi délicate, est précieux; nous n'en connoissons pas qui exige plus de dispositions naturelles: aussi nous sommes bien convaincus que, lorsque les savans auront fait, pour la perfection de cette méthode, tout ce que leur génie pourra leur inspirer, le talent de l'observateur aura encore beaucoup à y ajouter.
\end{quote}
\noindent Translation and typesetting produced by AI.
\clearpage
